%% Graphe orienté du fil conducteur du chapitre : Séries de Fourier ->
%% Transformée de Fourier -> TFD, chaque noeud légendé par le type de
%% fonctions concerné. Le noeud intermédiaire "Discret <=> Périodique"
%% a été délibérément retiré : le récit est "on part du périodique,
%% on va vers le discret" -- la périodicité en fréquence est une
%% conséquence qui apparaît en cours de route (pour obtenir une TFD
%% finie), pas une étape à part entière sur ce graphe. Requiert les
%% styles beatbox/beatfilled/beatarr (tikzset dans le préambule de
%% slides.tex). Variante "plain" -- aucun noeud rempli. Consommé par
%% slides (beat "Rappel").
\begin{tikzpicture}[node distance=0.6cm]
    \node[beatbox] (a) {Séries de\\Fourier};
    \node[below=of a,font=\small, text width=2.4cm, align=center](a-label){Fonctions continues périodiques};
    \node[beatbox, right=of a] (b) {Transformée\\de Fourier};
    \node[below=of b,font=\small, text width=2.4cm, align=center](b-label){Fonctions continues};
    \node[beatbox, right=of b] (c) {Transformée de Fourier Discrète};
    \node[below=of c,font=\small, text width=2.4cm, align=center](c-label){Fonctions discrètes};
    \draw[beatarr] (a) -- (b);
    \draw[beatarr] (b) -- (c);
\end{tikzpicture}
